\documentclass[./测控三级项目.tex]{subfiles} % 使用 subfiles 文档类，指定主文档

\begin{document}
	\subsection{题目重述}
		\begin{compactenum}
			\item 完成课本例2.1、2.2、2.4的仿真测试。
			\item 对I/V转换电路的理论值(依据公式计算值)与仿真值进行比较，计算相对误差。
			\item 简述三个电路的典型应用。
		\end{compactenum}
	\subsection{题目分析与实施方案}
		\subsubsection{第一问}
			通过课本查询可知：\par
			\textbf{例题2.1}\qquad 用741运算放大器和$R=$1M$\Omega$实现图\ref{fig:基本I-V转换器原理图}，并求它的闭环参数；
			\begin{figure}[h]
				\centering
				\includegraphics[width=0.45\textwidth]{基本I-V转换器原理图}
				\caption{基本I-V转换器原理图}
				\label{fig:基本I-V转换器原理图}
			\end{figure}\par
			\textbf{例题2.2}\qquad 在图\ref{fig:高灵敏度I-V转换器原理图}的电路中给出合适的元件值以实现0.1V/nA的灵敏度；
			\begin{figure}[h]
				\centering
				\includegraphics[width=0.45\textwidth]{高灵敏度I-V转换器原理图}
				\caption{高灵敏度I-V转换器原理图}
				\label{fig:高灵敏度I-V转换器原理图}
			\end{figure}\par
			\textbf{例题2.2}\qquad 根据图\ref{fig:Howland电流泵原理图}中的Howland泵，假设一个满程的运算放大器（$\pm V_{sat}=\pm 9V$)，准备一个表格显示$v_L =0,1,2,3,4,5,-2,-4,-6$情况下所有的电压和电流，并给出电路运行情况的文字描述。这个泵的电压柔量是多少？
			\begin{figure}[h]
				\centering
				\includegraphics[width=0.45\textwidth]{Howland电流泵原理图}
				\caption{Howland电流泵原理图}
				\label{fig:Howland电流泵原理图}
			\end{figure}\par
			整理课本知识点可知。\textbf{例题2.1}，\textbf{2.2}，\textbf{2.4}分别用到了基本I/V转换器，高灵敏度转换器，以及接地负载转换器的相关内容，该题要求对这三个电路仿真，需要使用Multisim软件仿真对应电路。\textbf{例题2.1}需要使用741运算放大器以及1M$\Omega$的电阻，仿真得到对应输出下的闭环参数，\textbf{例题2.2}需要根据公式反推出对应电阻之间的关系，并合理假设电阻阻值，得到符合题目要求的匹配电阻值，\textbf{例题2.4}根据题中给的输入电压值，对电路进行仿真，得到轨对轨运放在不同输入下输出值变化情况。
		\subsubsection{第二问}
			第二问要求对I/V转换电路的理论值（依据公式计算值）与仿真值进行比较，计算相对误差，要在书上找到求理论值的公式。得出结果后与第一问的仿真值进行比较，再计算相对误差。
		\subsubsection{第三问}
			第三问要求简述三个电路的典型应用，了解这三个电路各自是什么电路。然后通过书籍，浏览器等方式搜索三种电路的典型应用。
		
	\subsection{理论分析与仿真}
		\subsubsection{第一问}
			\textbf{例题2.1}\qquad 依据图\ref{fig:基本I-V转换器原理图}中所示电路，在Multisim软件中搭建如图\ref{fig:基本I-V转换器仿真原理图}电路。
			\begin{figure}[H]
				\centering
				\includegraphics[width=0.35\textwidth]{./5.第五题示波器现象/基本I-V转换器仿真原理图}
				\caption{基本I-V转换器仿真原理图}
				\label{fig:基本I-V转换器仿真原理图}
			\end{figure}\par
			为了防止输出电压受到电源输入影响导致失真，同时也为了能够快速计算得到仿真的增益，将电流源的值设置在了1$\mu$A。\par
			运行后电压表示数如图\ref{fig:第五题第一问电压表读数}所示，为-929.727mV。计算得到仿真的基本I-V转换器闭环增益为$A=\frac{-0.929727V}{1\mu A}=-0.929727$V/$\mu$A。
			\begin{figure}[H]
				\centering
				\includegraphics[width=0.30\textwidth]{./5.第五题示波器现象/第一问电压表读数}
				\caption{基本I-V转换器仿真电压表读数}
				\label{fig:第五题第一问电压表读数}
			\end{figure}\par
			\textbf{例题2.2}\qquad 依据图\ref{fig:高灵敏度I-V转换器原理图}中所示电路，在Multisim软件中搭建如图\ref{fig:高灵敏度I-V转换器仿真原理图}电路。
			\begin{figure}[H]
				\centering
				\includegraphics[width=0.35\textwidth]{./5.第五题示波器现象/高灵敏度I-V转换器仿真原理图}
				\caption{高灵敏度I-V转换器仿真原理图}
				\label{fig:高灵敏度I-V转换器仿真原理图}
			\end{figure}\par
			为了防止输出电压受到电源输入影响导致失真，同时也为了能够快速计算得到仿真的增益，将电流源的值设置在了1nA。\par
			运行后电压表示数如图\ref{fig:第五题第二问电压表读数}所示，为-100.044mV。计算得到仿真的基本I-V转换器闭环增益为$A=\frac{-0.100044V}{1\mu A}=-0.100044$V/nA。
			\begin{figure}[H]
				\centering
				\includegraphics[width=0.30\textwidth]{./5.第五题示波器现象/第二问电压表读数}
				\caption{高灵敏度I-V转换器仿真电压表读数}
				\label{fig:第五题第二问电压表读数}
			\end{figure}\par
			\textbf{例题2.4}\qquad 依据图\ref{fig:Howland电流泵原理图}中所示电路，在Multisim软件中搭建如图\ref{fig:Howland电流泵仿真原理图}电路。
			\begin{figure}[H]
				\centering
				\includegraphics[width=0.35\textwidth]{./5.第五题示波器现象/Howland电流泵仿真原理图}
				\caption{Howland电流泵仿真原理图}
				\label{fig:Howland电流泵仿真原理图}
			\end{figure}\par
			根据\textbf{例题2.4}题目要求，对电流表XMM1的内阻修改，达到修改负载电阻的效果，以此来得到不同的负载电压$v_L =0,1,2,3,4,5,-2,-4,-6$几种情况进行测试。得到关键参数如表\ref{tab:不同负载电压下Howland电流泵仿真电压和电流相关参数}所示。
			\begin{table}[h!]
				\centering
				\caption{不同负载电压下Howland电流泵仿真电压和电流相关参数}
				\label{tab:不同负载电压下Howland电流泵仿真电压和电流相关参数}
				\begin{tabular}{|c|c|c|c|c|}
					\hline
					$v_L$ & $i_1$  & $v_{OA}$ & $i_2$ & $i_O$ \\
					\hline
					0V & 1mA & 2.217mV & 1.109$\mu$A & 1.001mA\\
					\hline
					1V & 499.503$\mu$A & 2.004V & 501.566$\mu$A & 1.001mA\\
					\hline
					2V & -1.212$\mu$A & 4.007V & 1.002mA & 1.001mA\\
					\hline
					3V & -501.064$\mu$A & 6.006V & 1.502mA & 1.001mA\\
					\hline
					4V & -1.003mA & 8.012V & 2.004mA & 1.001mA\\
					\hline
					5V & -1.107mA & 8.114V & 1.95mA & 842.881$\mu$A\\
					\hline
					-2V & 2.001mA & -4.001V & -999.539$\mu$A & 1.001mA\\
					\hline
					-4V & 3.037mA & -8.112V & -2.019mA & 1.019mA\\
					\hline
					-6V & 2.834mA & -8.115V & -2.223mA & 0.611mA\\
					\hline
				\end{tabular}
			\end{table}	\par
			
		\subsubsection{第二问}
			\textbf{例题2.1}\qquad 图\ref{fig:基本I-V转换器原理图}是一个基本I-V转换器原理图。\par
			由于差分输入电阻$r_{id}$较大，一般认为同相和反相输入端直接没有电流流经，直接忽略该电流。于是，可依据电路结构，写出同相输入端$v_P$和反相输入端$v_N$的电压公式
			\begin{align}
				v_P&=v_I \\
				\frac{v_N-v_O}{R}&=i_I
			\end{align}\par
			结合等式$u_O=A_{od}u_I=A_{od}(v_P-v_N)$，联立可得到
			\begin{align}
				v_O=A_{od}(-Ri_I-v_O)
			\end{align}\par
			定义实际放大倍数$A$为输出电压$v_O$和输入电流$i_I$之比，则有
			\begin{align}\label{eq:实际基本I-V转换器放大倍数}
				A=\frac{v_O}{i_I}=-\frac{A_{od}R}{1+A_{od}}
			\end{align}\par
			本题使用开环增益$A_{od}=200$V/mV的运算放大器741，和电阻$R=1$M$\Omega$。根据实际放大倍数计算公式式\ref{eq:实际基本I-V转换器放大倍数}，可以计算得到实际放大倍数为-0.999995V/$\mu$A，和仿真算出的-0.929727V/$\mu$A相比，误差约为7\%在可以接受的范围之内，考虑到仿真环境中没有完全忽略差分输入电阻$r_{id}$，因此会有差距是正常现象。\par
			\textbf{例题2.2}\qquad 图\ref{fig:高灵敏度I-V转换器原理图}是一个高灵敏度I-V转换器原理图。\par
			由于差分输入电阻$r_{id}$较大，一般认为同相和反相输入端直接没有电流流经，直接忽略该电流。于是，可依据电路结构，写出同相输入端$v_P$和反相输入端$v_N$的电压公式
			\begin{align}
				v_P&=0 \\
				i_I&=\frac{v_N-v_B}{R}\\
				\frac{v_N-v_B}{R}+\frac{0-v_B}{R_1}&=\frac{v_B-v_O}{R_2}
			\end{align}\par
			结合等式$u_O=A_{od}u_I=A_{od}(v_P-v_N)$，联立可得到
			\begin{align}
				v_O=A_{od}(-\frac{R(RR_1+RR_2+R_1R_2)}{RR_1+RR_2}i_I-\frac{RR_1}{RR_1+RR_2}v_O)
			\end{align}\par
			定义实际放大倍数$A"$为输出电压$v_O$和输入电流$i_I$之比，则有
			\begin{align}\label{eq:实际高灵敏度I-V转换器放大倍数}
				A"=\frac{v_O}{i_I}=-\frac{A_{od}R(RR_1+RR_2+R_1R_2)}{RR_1+RR_2+A_{od}RR_1}
			\end{align}\par
			本题使用开环增益$A_{od}=+\infty$V/mV的运算放大器741，和电阻$R=1$M$\Omega$、$R_1=1$k$\Omega$、$R_2=99$k$\Omega$。根据实际放大倍数计算公式式\ref{eq:实际基本I-V转换器放大倍数}，可以计算得到实际放大倍数为-0.1V/nA，和仿真算出的-0.100044V/nA相比，误差约为0.044\%在可以接受的范围之内，考虑到仿真环境中没有完全忽略差分输入电阻$r_{id}$，因此会有差距是正常现象。\par
			
			\textbf{例题2.4}\qquad 图\ref{fig:Howland电流泵原理图}是一个Howland电流泵原理图。\par
			由于差分输入电阻$r_{id}$较大，一般认为同相和反相输入端直接没有电流流经，直接忽略该电流。于是，可依据电路结构，写出同相输入端$v_P$和反相输入端$v_N$的电压公式。
			\begin{align}
				\frac{R_3}{R_3+R_4}v_O&=v_N	\label{eq:Howland泵负输入端}\\
				\frac{v_P}{R_L}&=i_O		\label{eq:Howland泵负载}\\
				\frac{v_I-v_P}{R_1}+\frac{v_O-v_P}{R_2}&=i_O\label{eq:Howland泵正输入端}
			\end{align}\par
			结合等式$u_O=A_{od}u_I=A_{od}(v_P-v_N)$、式\ref{eq:Howland泵负输入端}和式\ref{eq:Howland泵负载}，联立可得到
			\begin{align}
				v_O=A_{od}(R_Li_O-\frac{R_3}{R_3+R_4}v_O)
				\label{eq:Howland泵运放公式}
			\end{align}\par
			将式\ref{eq:Howland泵负载}和式\ref{eq:Howland泵运放公式}代入式\ref{eq:Howland泵正输入端}中，并定义实际放大倍数$A$为输出电流$i_O$和输入电压$v_I$之比，则整理可得：
			\begin{align}
				A=\frac{i_O}{v_I}=-\frac{(R_3+R_4+A_{od}R_3)R_2}{A_{od}(R_LR_1R_4-R_LR_2R_3-R_1R_2R_3)-(R_LR_1+R_LR_2+R_1R_2)(R_3+R_4)}
			\end{align}\par
			对于本题仿真，除$R_L$以外所有电阻采用$2$k$\Omega$，即$R_1=R_2=R_3=R_4=2$k$\Omega$。
			\begin{table}[h!]
				\centering
				\caption{不同负载电压下Howland电流泵电压和电流相关参数理论仿真对比}
				\label{tab:不同负载电压下Howland电流泵电压和电流相关参数理论仿真对比}
				\begin{tabular}{|c|c|c|c|c|c|c|c|c|c|c|c|c|}
					\hline
					\multirow{2}{*}{$v_L$} & \multicolumn{3}{c|}{$i_1$} & \multicolumn{3}{c|}{$v_{OA}$}  & \multicolumn{3}{c|}{$i_2$} &\multicolumn{3}{c|}{$i_O$}  \\
					\cline{2-13}
					 & 计算 & 实测 & 误差 & 计算 & 实测 & 误差 & 计算 & 实测 & 误差 & 计算 & 实测 & 误差\\
					\hline
					0 &  1.0 &  1.000 & 0.00\% &0 &0.002 & /	  &0.0 &0.001 & / & 1 & 1.001 & 0.10\%\\
					\hline
					1 &  0.5 &  0.499 & 0.10\% &2 &2.004 & 0.20\% &0.5 &0.501 & 0.31\% & 1 & 1.001 &  0.10\%\\
					\hline
					2 &  0.0 & -0.001 & /      &4 &4.007 & 0.18\% &1 &1.002 & 0.20\% & 1 & 1.001 & 0.10\%\\
					\hline
					3 & -0.5 & -0.501 & 0.21\% &6 &6.006 & 0.10\% &1.5 &1.502 & 0.13\% & 1 & 1.001 & 0.10\%\\
					\hline
					4 & -1.0 & -1.003 & 0.30\% &8 &8.012 & 0.15\% &2.0 &2.004 & 0.20\% & 1 & 1.001 & 0.10\%\\
					\hline
					5 & -1.5 & -1.107 & 26.20\% &9 &8.114 & 9.84\% &2.0 &1.95 & 2.50\% & 0.5 & 0.842 & 68.57\%\\
					\hline
					-2 &2.0 &  2.001 & 0.05\% &-4 &-4.001 & 0.03\% &-1 &-0.999 & 0.05\% & 1 & 1.001 & 0.10\%\\
					\hline
					-4 &3.0 &  3.037 & 1.23\% &-8 &-8.112 & 1.40\% &-2 &-2.019 & 0.95\% & 1 & 1.019 & 1.90\%\\
					\hline
					-6 &4.0 &  2.834 & -29.15\% &-9 &-8.115 & 9.83\% &-1.5 &-2.223 & 48.20\% & 2.5 & 0.611 & 75.56\%\\
					\hline
				\end{tabular}
			\end{table}	\par
			对于电压输出不超过运放电源电压的时候，Howland泵均能正常工作，并给出精度较高的电流源电流输出；对于运放电压输出超过运放电源电压的时候，输出电压受到限制，导致了电流源的电流输出失真。
			
		\subsubsection{第三问}
			基本I/V转换器的典型应用有：
			\begin{enumerate}
				\item 对电流控制负载的精确控制。例如，在电机控制系统中，通过调整电流来控制电机的转速和转矩。
				\item 接在传感器测量输出的电流信号之后，进行后续电信号转化运算，方便进行模拟信号处理和采集（如使用ADC采集成数字信号）。
				\item 作为功率放大器，将输入的电压信号放大后输出到负载上。这种应用在音频设备、无线电通信等领域中较为常见。
				\item 构建模拟前端，接收和处理来自天线和传感器的信号。例如，在电离层遥感探测和全球航海通信中，电流/电压转换器用于接收和处理甚低频信号
			\end{enumerate}
			高敏感I/V转换器的典型应用有：
			\begin{enumerate}
			\item 作为高输入阻抗和高灵敏度运算放大器的应用实例，用于测量束流强度等弱电流。
			\item 应用于电离层遥感探测与通信在电离层遥感探测和地球空间环境地基监测中，通过分析、仿真电流/电压转换器构建模拟前端的灵敏度和频率响应，能有效接收自然界雷电辐射和人工台站发射的甚低频信号
			\item 应用于在智能扫地机器人的激光测距模块中，需要将激光管输出的电流信号转换为电压信号，并进行放大处理，最后传送给MCU处理。
			\item 应用于微弱信号检测和高增益放大中，也有广泛应用。
			\end{enumerate}
			Howland泵的典型应用有：
			\begin{enumerate}
				\item 在电阻层析成像技术中被广泛使用，作为电压转换电流源的基本形式。其稳定性和精确性对于成像质量至关重要
				\item 在信号传输过程中，特别是需要远距离传输时，可以将电压信号转换为电流信号，利用电流信号对噪声不敏感的特性，提高信号传输的抗干扰能力和稳定性。
				\item 在一些复杂的电源系统中，被用作电压控制电流源（VCCS），用于选择和控制电压源和电流源。这使得系统能够根据不同的应用需求，灵活地切换电源类型，提高系统的通用性和适应性。
				\item 被用作电流源或电流放大器。通过调整电路中的电阻和运放参数，可以实现对输出电流的精确控制。
			\end{enumerate}
			
			
	
	
	\subsection{结论与收获}
		\subsubsection{结论}
		经验证，所设计三种电路基本满足题目要求，仿真值与计算值相比，略有误差但在允许范围之内。
		\subsubsection{收获}
		\textbf{游寅龙}：\par
		理论分析层面：学习了这些电路的基本原理及其背后的数学模型，对例题的理论推导关系进行了计算，成功推导出正确的关系。\par
		动手实践层面：参与了仿真电路的设计和检验，记录了所得数据，能够将理论用于解决实际问题。\par
		精神层面：体会到了坚持与耐心的价值，面对电路推导时要时刻注意电流方向等细节，逐步分析解决问题，培养了不放弃的精神。\par
		\textbf{郭嘉媛}：\par
		理论分析层面：系统的学习了三种电路的工作原理，对三种电路各自的工作方式有了更具体地理解。\par
		动手实践方面：利用推导的关系式完成了Multsim的电路仿真，能够根据错误地实验结果反推错误的环节，培养了自我纠错能力。\par
		精神层面：在纠错时，需要足够的细心，不放过每一个可能的错误，培养了严谨的科学实验精神。\par
		\textbf{郑嘉宁}：\par
		理论分析层面：将课本例题与仿真电路进行对比，还对三种电路的应用进行查找，拓宽了视野，对三种电路的理解也更深一步。\par
		动手实践层面：对所得数据进行了整理，分析了计算值与仿真值的相对误差，增强了数据处理的能力。\par
		精神层面：在分析数据的过程中积极与小组成员进行交流，得到了小组成员的许多帮助，增强了团队意识。\par
		
		
	
	
\end{document}